% Part: sets-functions-relations
% Chapter: functions
% Section: inverses

\documentclass[../../../include/open-logic-section]{subfiles}

\begin{document}

\olfileid{sfr}{fun}{inv}
\olsection{অপেক্ষকের বিপরীত}

\begin{explain}
আমরা অপেক্ষককে চিত্রণ হিসেবে ভাবি। তাই স্বাভাবিকভাবেই প্রশ্ন ওঠে, এই চিত্রণের কাজ ``উল্টে দেওয়া'' যায় কি না। যেমন, উত্তরসূরি অপেক্ষক $f(x) = x + 1$-এর কাজ এই অর্থে উল্টে দেওয়া যায় যে অপেক্ষক $g(y) = y - 1$, $f$-এর করা কাজটিকে ``পূর্বাবস্থায় ফিরিয়ে দেয়''।

তবে সাবধান হতে হবে। যদিও~$g$-এর সংজ্ঞাটি একটি অপেক্ষক $\Int \to \Int$ নির্ধারণ করে, এটি একটি \emph{অপেক্ষক} $\Nat
\to \Nat$ নির্ধারণ করে না, কারণ $g(0) \notin \Nat$। সুতরাং সরল ক্ষেত্রেও একটি অপেক্ষকের কাজ উল্টে দেওয়া যায় কি না, তা একেবারে স্পষ্ট নয়; উত্তরটি সংজ্ঞাক্ষেত্র ও সহসংজ্ঞাক্ষেত্রের উপর নির্ভর করতে পারে।

অপেক্ষকের বিপরীতের ধারণা দিয়ে বিষয়টি আরও নির্দিষ্ট করা যায়।
\end{explain}

\begin{defn}
একটি অপেক্ষক $g \colon B \to A$-কে অপেক্ষক $f
\colon A \to B$-এর \emph{বিপরীত} বলা হয়, যদি সকল $x \in A$ ও $y \in B$-এর জন্য $f(g(y)) = y$ এবং $g(f(x)) = x$ হয়।
\end{defn}

$f$-এর বিপরীত~$g$ থাকলে, আমরা প্রায়ই~$g$-এর পরিবর্তে $f^{-1}$ লিখি।

\begin{explain}
এখন আমরা নির্ধারণ করব, কখন অপেক্ষকের বিপরীত থাকে। $f\colon A \to B$-এর বিপরীতের একটি সম্ভাব্য প্রার্থী হলো নিচের নিয়মে ``সংজ্ঞায়িত'' $g\colon B \to A$:
\[
g(y) = \text{``সেই'' $x$, যার জন্য $f(x) = y$।}
\]
কিন্তু ``সংজ্ঞায়িত'' (এবং ``সেই'') কথার চারপাশের উদ্ধৃতিচিহ্ন ইঙ্গিত দেয় যে এটি সত্যিকার সংজ্ঞা নয়। অন্তত সম্পূর্ণ সাধারণভাবে এটি সব সময় কাজ করবে না। কারণ এই সংজ্ঞা দিয়ে একটি অপেক্ষক নির্দিষ্ট করতে হলে এমন এক এবং কেবল এক~$x$ থাকতে হবে, যাতে $f(x) =
y$ হয়—$g$-এর আউটপুট অনন্যভাবে নির্দিষ্ট হতে হবে। তদুপরি, প্রত্যেক $y \in B$-এর জন্য তা নির্দিষ্ট হতে হবে। যদি এমন $x_1$ ও $x_2 \in
A$ থাকে, যাতে $x_1 \neq x_2$ অথচ $f(x_1) = f(x_2)$, তবে $y = f(x_1) = f(x_2)$-এর জন্য $g(y)$ অনন্যভাবে নির্দিষ্ট হবে না। আর যদি $f(x) = y$ হয় এমন কোনো~$x$-ই না থাকে, তবে $g(y)$ আদৌ নির্দিষ্ট হবে না। অন্যভাবে বললে, $g$-কে সংজ্ঞায়িত করতে হলে $f$-কে একই সঙ্গে !!{injective} ও !!{surjective} হতে হবে।

ধীরে এগোনো যাক। প্রশ্নটিকে দুটি ভাগে ভাগ করব: একটি অপেক্ষক~$f\colon A \to B$ দেওয়া থাকলে, কখন এমন একটি অপেক্ষক $g\colon B \to A$ থাকে, যাতে $g(f(x)) = x$ হয়? এই ধরনের $g$, $f$-এর কাজকে ``পূর্বাবস্থায় ফিরিয়ে দেয়''; একে~$f$-এর \emph{বাম বিপরীত} বলা হয়। দ্বিতীয়ত, কখন এমন একটি অপেক্ষক $h\colon B \to A$ থাকে, যাতে $f(h(y)) = y$ হয়? এই ধরনের $h$-কে~$f$-এর \emph{ডান বিপরীত} বলা হয়—এক্ষেত্রে $f$, $h$-এর কাজকে ``পূর্বাবস্থায় ফিরিয়ে দেয়''।
\end{explain}

\begin{prop}
সংজ্ঞাক্ষেত্রটি অশূন্য এবং $f\colon A \to B$ যদি !!{injective} হয়, তবে~$f$-এর এমন একটি \emph{বাম বিপরীত}~$g\colon B \to A$ আছে, যাতে সকল $x
\in A$-এর জন্য $g(f(x)) = x$ হয়।
% OLFUN-001 TARGET-CORRECTION-BEGIN
\par\smallskip\noindent\textnormal{\small\textbf{উৎস-সংশোধন (OLFUN-001):} হিমায়িত ইংরেজি প্রতিজ্ঞায় $A$ অশূন্য হওয়ার শর্তটি নেই, যদিও প্রমাণে $A$ থেকে একটি $a$ বেছে নেওয়া হয়েছে। $A=\emptyset$ ও $B=\{0\}$ হলে শূন্য অপেক্ষকটি একৈক, কিন্তু $B\to A$ অপেক্ষক নেই। সঠিক সাধারণ শর্ত: একৈক $f\colon A\to B$-এর বাম বিপরীত থাকে যদি এবং কেবল যদি $A$ অশূন্য অথবা $B$ শূন্য।} % OLFUN-001 TARGET-CORRECTION-END
\end{prop}

\begin{proof}
ধরি, সংজ্ঞাক্ষেত্রটি অশূন্য এবং $f\colon A \to B$ !!{injective}। একটি $y \in B$ বিবেচনা করো। যদি $y \in \ran{f}$ হয়, তবে এমন একটি $x \in A$ আছে, যাতে $f(x) = y$ হয়। যেহেতু $f$ !!{injective}, তাই এমন~$x \in A$ একটিই আছে। তখন আমরা $g(y) = x$ সংজ্ঞায়িত করতে পারি, অর্থাৎ $g(y)$ হলো ``সেই'' $x \in A$, যার জন্য $f(x)
= y$। যদি $y \notin \ran{f}$ হয়, তবে তাকে যেকোনো~$a \in A$-তে পাঠাতে পারি। তাই একটি $a \in A$ বেছে নিয়ে নিচের নিয়মে $g\colon B \to A$ সংজ্ঞায়িত করতে পারি:
\[
g(y) = \begin{cases}
    x & \text{যদি $f(x) = y$ হয়}\\
    a & \text{যদি $y \notin \ran{f}$ হয়।}
\end{cases}
\]
এটি সকল $y \in B$-এর জন্য সংজ্ঞায়িত, কারণ এই ধরনের প্রত্যেক $y \in \ran{f}$-এর জন্য ঠিক একটি $x \in A$ আছে, যাতে $f(x) = y$ হয়। সংজ্ঞা অনুসারে, $y = f(x)$ হলে $g(y) = x$, অর্থাৎ $g(f(x)) = x$।
\end{proof}

\begin{prob}
দেখাও যে $f\colon A \to B$-এর একটি বাম বিপরীত~$g$ থাকলে $f$ !!{injective} হয়।
\end{prob}

\begin{prop}
    $f\colon A \to B$ যদি !!{surjective} হয়, তবে~$f$-এর এমন একটি \emph{ডান বিপরীত}~$h\colon B \to A$ আছে, যাতে সকল~$y \in B$-এর জন্য $f(h(y)) =
    y$ হয়।
\end{prop}

\begin{proof}
ধরি, $f\colon A \to B$ !!{surjective}। একটি $y \in
B$ বিবেচনা করো। যেহেতু $f$ !!{surjective}, তাই এমন একটি $x_y \in A$ আছে, যাতে $f(x_y)
= y$। তখন আমরা $h(y) = x_y$ সংজ্ঞায়িত করতে পারি, অর্থাৎ প্রত্যেক $y \in B$-এর জন্য এমন কোনো $x \in A$ বেছে নিই, যাতে $f(x) = y$ হয়; যেহেতু $f$ !!{surjective}, তাই বেছে নেওয়ার মতো অন্তত একটি উপাদান সব সময় আছে।\footnote{যেহেতু $f$ !!{surjective}, তাই প্রত্যেক~$y \in B$-এর জন্য সেট $\Setabs{x}{f(x) = y}$ অশূন্য। আমাদের~$h$-এর সংজ্ঞার জন্য এই সেটগুলির প্রত্যেকটি থেকে একটি করে $x$ বেছে নিতে হয়। এটি যে সব সময় সম্ভব, তা আসলে স্পষ্ট নয়—এই নির্বাচনগুলি করা সম্ভব বলে একটি স্বতঃসিদ্ধের মাধ্যমে ধরে নেওয়া হয়। অন্যভাবে বললে, এই প্রতিজ্ঞাটি তথাকথিত নির্বাচন স্বতঃসিদ্ধ ধরে নেয়; বিষয়টি আমরা \oliflabeldef{sth:choice::chap}{\olref[sth][choice][]{chap}-এ আবার আলোচনা করব}{এখানে বিস্তারিত আলোচনা করব না}। তবে অনেক নির্দিষ্ট ক্ষেত্রে, যেমন $A = \Nat$ হলে বা সেটি সসীম হলে, অথবা $f$ !!{bijective} হলে, নির্বাচন স্বতঃসিদ্ধের প্রয়োজন হয় না। (বিশেষত $f$ !!{bijective} হলে, প্রত্যেক $y \in B$-এর জন্য সেট $\Setabs{x}{f(x) = y}$-এ ঠিক একটি !!{element} থাকে, ফলে কোনো নির্বাচন করতে হয় না।)} সংজ্ঞা অনুসারে, $x = h(y)$ হলে $f(x) = y$, অর্থাৎ যেকোনো $y \in B$-এর জন্য $f(h(y)) = y$।
\end{proof}

\begin{prob}
দেখাও যে $f\colon A \to B$-এর একটি ডান বিপরীত~$h$ থাকলে $f$ !!{surjective} হয়।
\end{prob}

\begin{explain}
  আগের প্রমাণের ধারণাগুলি একত্র করলে এখন পাওয়া যায় যে প্রত্যেক !!{bijection}-এর বিপরীত আছে, অর্থাৎ এমন একটি অপেক্ষক আছে, যা একই সঙ্গে~$f$-এর বাম ও ডান বিপরীত।
\end{explain}

\begin{prop}\ollabel{prop:bijection-inverse}
$f\colon A \to B$ যদি !!{bijective} হয়, তবে এমন একটি অপেক্ষক~$f^{-1}\colon B \to A$ আছে, যাতে সকল $x \in A$-এর জন্য $f^{-1}(f(x)) = x$ এবং সকল $y \in B$-এর জন্য $f(f^{-1}(y)) = y$ হয়।
\end{prop}

\begin{proof}
অনুশীলনী।
\end{proof}

\begin{prob}
\olref[sfr][fun][inv]{prop:bijection-inverse} প্রমাণ করো। তোমাকে~$f^{-1}$ সংজ্ঞায়িত করতে হবে, সেটি একটি অপেক্ষক দেখাতে হবে, এবং সেটি~$f$-এর বিপরীত দেখাতে হবে, অর্থাৎ সকল $x \in A$ ও $y \in B$-এর জন্য $f^{-1}(f(x)) = x$ এবং $f(f^{-1}(y)) = y$ হয় দেখাতে হবে।
\end{prob}

\begin{explain}
বিপরীত পাওয়ার কিছুটা বেশি সাধারণ একটি পদ্ধতি আছে। \olref[kin]{sec}-এ আমরা দেখেছি, প্রত্যেক অপেক্ষক $f$, সকল $x \in A$-এর জন্য $f'(x) = f(x)$ রেখে !!a{surjection} $f'
\colon A \to \ran{f}$ নির্ধারণ করে। স্পষ্টতই $f$ !!{injective} হলে $f'$ !!{bijective}, ফলে \olref{prop:bijection-inverse} অনুসারে তার একটি অনন্য বিপরীত আছে। সংকেতের ব্যবহারে অতি সামান্য শিথিলতা মেনে নিয়ে, আমরা কখনো কখনো $f'$-এর বিপরীতকে শুধু ``$f$-এর বিপরীত'' বলি।
\end{explain}

\begin{prop}\ollabel{prop:left-right}%
  দেখাও যে $f\colon A \to B$-এর একটি বাম বিপরীত~$g$ এবং একটি ডান বিপরীত~$h$ থাকলে $h = g$ হয়।
\end{prop}

\begin{proof}
  অনুশীলনী।
\end{proof}

\begin{prob}
  \olref[sfr][fun][inv]{prop:left-right} প্রমাণ করো।
\end{prob}

\begin{prop}\ollabel{prop:inverse-unique}
প্রত্যেক অপেক্ষক~$f$-এর বড়জোর একটি বিপরীত থাকে।
\end{prop}

\begin{proof}
  ধরি, $g$ ও $h$ উভয়েই~$f$-এর বিপরীত। তাহলে বিশেষত $g$ হলো~$f$-এর একটি বাম বিপরীত এবং $h$ একটি ডান বিপরীত। \olref{prop:left-right} অনুসারে $g = h$।
\end{proof}

\end{document}
